\documentclass[letterpaper,twoside,openright]{book}
%\usepackage[left=4cm,right=4cm,top=4cm,bottom=4cm,letterpaper]{geometry}
\usepackage[spanish]{babel}
\usepackage[utf8]{inputenc}
\usepackage{graphicx}
\usepackage{mathptmx}

%\author{Héctor E. Gómez Morales}
%\title{Reporte Trabajo Profesional: Sistema PEAT (Progressive Energy
%  Audit Tool)}

\begin{document}

\thispagestyle{empty}

\frontmatter

\begin{minipage}{.3\textwidth}
  \flushleft
  \center{\includegraphics[scale=.09]{imagenes/unam.pdf}}

  \vspace{20pt}

  \center{
    \rule{.5pt}{.6\textheight}
    \hspace{7pt}
    \rule{2pt}{.6\textheight}
    \hspace{7pt}
    \rule{.5pt}{.6\textheight}
  } \\

  \center{\includegraphics[scale=.22]{imagenes/ciencias.pdf}}
\end{minipage}
\begin{minipage}{.7\textwidth}
\flushright

\center{

  \center{
    \LARGE{U}\large{NIVERSIDAD} \LARGE{N}\large{ACIONAL} 
    \LARGE{A}\large{UTÓNOMA} \\[10pt]
    \large{DE} 
    \LARGE{M}\large{ÉXICO} 
  } \\
  \rule{\textwidth}{2pt}
  \\
  \hrulefill\\[1cm]
  
  \LARGE{F}\large{ACULTAD DE } \LARGE{C}\large{IENCIAS}\\[2cm]

  \large{
Diseño e implementación del sistema PEAT (Progressive Energy Audit Tool)  }\\[2cm]

  \huge{
REPORTE DE TRABAJO PROFESIONAL  }\\[1cm]

  \large{QUE PARA OBTENER EL TÍTULO DE:}\\[1cm]

  \large{
Licenciado en Ciencias de la Computación  }\\[1cm]

  \large{PRESENTA:}\\[1cm]

  \large{
Héctor Enrique Gómez Morales  }\\[1cm]

  \large{
TUTORA  }\\[1cm]

  \large{
Karla Ramírez Pulido  }
}

\end{minipage}

%\maketitle
\tableofcontents

\chapter{Introducción}
\section*{Progressive Energy Audit Tool (PEAT)}
Este trabajo aborda el diseño e implementación del sistema \textit{Progressive
  Energy Audit Tool}\ (PEAT), el cual es un sistema web que permite a los usuarios
de pequeñas y medianas empresas (PyMES) identificar y monitorear sus gastos
de energía eléctrica, con el fin de obtener un mayor control en los costos de ésta.

Este sistema fue desarrollado siendo parte de la compañía
\textbf{Software Next Door} en el cargo de \textbf{Senior Software Engineer}.

Pacific Gas and Electric Company (PG\&E) es una compañía proveedora de gas natural
y electricidad, una de las más grandes compañías de Estados Unidos con sede en
San Francisco, California. El sistema PEAT es el resultado de una licitación
iniciada por PG\&E para el desarrollo de un sistema enfocado en usuarios PyMES.

Su desarrollo es de importancia crítica tanto para PG\&E, como para el gobierno de
California puesto que su correcto funcionamiento es un requisito en una nueva
ley de facturación de energía eléctrica desde el 2013. La nueva ley buscaba
lograr distribuir la carga de la red eléctrica, siendo una de las formas
principales para lograr esto el desincentivar el uso de la red
eléctrica en horas pico al así darle la facultad a las compañías como PG\&E
de cobrar tasas mucho más altas en estas horas.

PG\&E había realizado el despliegue de medidores inteligentes en la
mayor parte de sus clientes, por lo que se tenía acceso a información
muy detallada del consumo de energía de éstos.

Actualmente para darle el mayor valor e información a las empresas es necesario
el obtener más datos de su entorno de operación; por ejemplo el número
de edificios asociados a la cuenta, rubro de la empresa, número
de empleados, etcétera. De esta forma, entre mas detalles se pueda captar sobre
la empresa el sistema tendrá más facilidad en obtener un desglose más robusto,
completo y útil de sus consumos de energía.

\mainmatter

\chapter{Situación actual}

\section{Contexto}
De 2000 a 2002 se tuvo una crisis energética en el estado de California en los
Estados Unidos de América, esta crisis provocó que el gobierno del estado empezara a
tomar decisiones de largo plazo en su estrategia energética. En los años
subsecuentes se aprobaron nuevas leyes para incentivar el uso de energía renovable,
la producción de energía, etcétera.

Para el 2010 se aprobó una nueva legislación que permitía un nuevo esquema
de cobro a las empresas generadoras de electricidad. Este nuevo esquema
llamado \textit{hora de consumo} (\textit{time-of-use TOU}) define que las
tasas de cobro varíen en función de la temporada y la hora del día en que
se usa la energía eléctrica, en contraste con las tasas fijas tradicionales.

En un esquema \texttt{TOU} se definen tres tarifas y sus zonas de tiempo asociadas:

\begin{itemize}
\item  Demanda baja: es la tarifa mas baja, se aplica durante las horas de la
  mañana, en la noche y los fines de semana.
\item Demanda mediana: es la tarifa intermedia, se aplica durante las horas de
  10 am a 1pm y de 7 pm a 9 pm.
\item Demanda Alta: es la tarifa mas cara, se aplica durante las horas de
  1 pm a 7 pm.
\end{itemize}

El objetivo de un esquema \texttt{TOU} es el distribuir la carga de la red
eléctrica, alentando a los usuarios a cambiar su consumo durante periodos
de demanda alta a periodos de demanda baja, al bajar la demanda en los periodos
de demanda alta provoca que se reduzcan los costos de generación de energía de
forma general.

Un esquema \texttt{TOU} es solo posible con la llegada de los medidores
inteligentes, a diferencia de los antiguos medidores analógicos los medidores
inteligentes son capaces de medir el consumo de electricidad de forma instantánea,
de distinguir y facturar el consumo de electricidad según el momento en que se
esta realizando el consumo. Desde inicios del 2010 Pacific Gas and
Electric Company (PG\&E) empezó la instalación de medidores inteligentes
en su territorio, por lo que para inicios del 2012 ya casi había
completado la transición de los medidores.

Como fecha de arranque de la nueva legislación se puso el mes de noviembre del 2012,
las empresas como PG\&E podían iniciar la transición a un esquema \texttt{TOU}
de sus usuarios PyMES. Una parte vital para permitir esta transición es que
PG\&E brindara herramientas para que el usuario final pudiera analizar sus
gastos de energía y tomar decisiones sobre su consumo de energía según las
nuevas tarifas y su historial de consumo.

En esta situación es que nace el sistema \textit{Progressive
  Energy Audit Tool} \ (PEAT) para solventar la necesidad de información
del usuario final sobre sus gastos de energía en una forma detallada.
La licitación para la implementación del sistema \texttt{PEAT} fue puesta
a concurso por parte PG\&E siendo la compañía C3 Energy la ganadora de dicha
licitación  ya que contaba con la infraestructura necesaria para el
procesamiento de una gran cantidad de datos de consumo de energía,
tiene un sistema de monitoreo de consumo de energía enfocada a empresas de nivel
multinacional. El reto era pasar de un sistema y procesos diseñados
para una docena de clientes, a un sistema que diera servicio a
cientos de miles de clientes PyMES.

\section{Objetivo general}

El objetivo general del sistema \texttt{PEAT} es dar la mayor utilidad posible
con la menor información disponible, fomentando en el usuario final el
compartir más información sobre su empresa, obteniendo un mejor control
acerca de su consumo energético.

El sistema debe permitir el ingreso progresivo de información, por medio de una
serie de preguntas especificas al usuario, dando mejores recomendaciones para
bajar su consumo energético conforme el sistema obtiene mas información.

El sistema debe permitir el monitoreo y revisión del consumo energético de forma
detallada, ya sea en horas, días, meses y años.

\section{Objetivos secundarios}

Los objetivos secundarios que apoyan al objetivo general de este trabajo son:
\begin{itemize}
\item Implementación de una interfaz que permita la obtención
  de información del usuario de una forma eficaz y sencilla.
\item Dar información útil aunque el usuario solo proporcione el
  mínimo de información sobre su empresa.
\item Proporcionar recomendaciones para disminuir sus
  gastos en energía con base en el consumo e información proporcionada
  hasta el momento.
\item Autentificar a los usuarios mediante el uso de credenciales de acceso
  obtenidas en el portal web de PG\&E.
\item Soportar por lo menos a mil usuarios concurrentes.
\item Implementar la infraestructura para el despliegue continuo de la
  aplicación, permitiendo una retroalimentación continua sobre el
  funcionamiento del sistema.
\end{itemize}

\section{Arquitectura previa}
En C3 Energy ya se contaba con toda una infraestructura para el
procesamiento y análisis en tiempo real de datos de consumo eléctrico
y de gas.

En grandes rasgos el sistema consta de cuatro capas una capa de
almacenamiento de información que era de referencia o datos ya
procesados, otra capa en las que se hacia el análisis de los datos
y que también tomaba el rol de capa de cache. La tercera capa era el
servidor que servía el API para permitía acceder y realizar análisis
sobre las capas de datos anteriores y finalmente la cuarta capa que
era el cliente que realizaba llamadas al API para presentar la
información al usuario final.

\subsection{Bases de datos}
Las dos primeras capas almacenan y analizan la información obtenida
de un gran numero de bases de datos externas relacionadas con temas
de energía. Todos estos datos son concentrados en una base de datos
central en la que se usaba Oracle Database 10g para este rol.

También se contaba con un cluster de Apache Cassandra el cual era usado
para realizar el análisis en tiempo real de los datos, además de servir
como capa de cache para los datos mas solicitados en el sistema.

\subsection{Servidor y API}
Sobre la anterior bases de datos se tenia un API que que permita el
hacer consultas y operaciones sobre los datos contenidos y/o analizados
en este sistema. Este API estaba implementado en Javascript bajo Rhino
que corre bajo la maquina virtual de Java (JVM), esto daba la ventaja
que se podía realizar la implementación de rutinas criticas en un
lenguaje con mejor rendimiento como Java. De esta forma se tenia que el
núcleo del servidor estaba implementado en Java con el resto del API
implementado en Javascript.

\subsection{Aplicaciones}
Las aplicaciones ya existentes para empresas multinacionales eran del tipo
SPA (Single Page Application) que son aplicaciones web que tienen el fin
de dar una experiencia de usuario similar a la de una aplicación de escritorio.
Esto era implementado usando un framework hecho dentro de la compañía
en Javascript. Se tenia una gran jerarquía de objetos que era
compartida tanto en el servidor como en el cliente, así se tenia que
había una sola fuente de la descripción de los objetos del sistema.

Se hacia un uso extenso de la biblioteca Ext JS que es un framework
en Javascript para construir aplicaciones del tipo SPA, se hacia un uso particular
a su facultad para rápidamente generar todo tipo de gráficas para mostrar
información al usuario.

Entonces tenemos que tanto en el frontend como en el backend se hacia
uso de Javascript principalmente con un núcleo de Java para las partes
que requerían un rendimiento óptimo. Esto tenia varias ventajas:

\begin{itemize}
\item Reducir al mínimo el cambio de contexto entre lenguajes entre
  el frontend y el backend.
\item Reducir sustancialmente la cantidad de código redundante al no
  tenerse que reimplementar la jerarquía de objetos, que era pasada
  por el servidor en forma serializada por JSON y deserializada por el
  cliente
\item Había una única representación de una clase por ejemplo la
  clase de edificios era una clase de Javascript que se usaba para
  tanto en el servidor como en el cliente.
\end{itemize}

\section{Propuesta de desarrollo}


\end{document}
